\section{최대공약수와 베주 항등식}
\label{sec:polynomial-gcd}

\begin{learninggoal}
$K[X]$에서 유클리드 알고리즘으로 최대공약수를 계산하고, 베주 항등식과 몫환에서의 역원 계산에 활용한다.
\end{learninggoal}

체 $K$에 대해 $K[X]$는 차수를 유클리드 함수로 갖는 유클리드 정역이다. 따라서 모든 아이디얼이 주아이디얼이고, 두 다항식의 최대공약수를 나눗셈을 반복하여 계산할 수 있다.

\begin{definition}
$f,g\in K[X]$가 둘 다 영이 아닐 때 일계수다항식 $d$가
\begin{enumerate}[label=(\roman*)]
  \item $d\mid f$이고 $d\mid g$이며,
  \item $h\mid f$와 $h\mid g$이면 $h\mid d$
\end{enumerate}
를 만족하면 $d$를 $f,g$의 최대공약수라 하고 $d=\gcd(f,g)$로 쓴다.
\end{definition}

일계수 조건은 단위원배의 모호성을 제거한다. $K[X]$의 단위원은 $K^\times$이므로 최대공약수는 상수배를 제외하면 유일하다.

\begin{theorem}[유클리드 알고리즘]
$f,g\in K[X]$, $g\ne0$라 하자. 연속적인 나눗셈
\[
\begin{aligned}
 f&=q_1g+r_1, &\deg r_1&<\deg g,\\
 g&=q_2r_1+r_2, &\deg r_2&<\deg r_1,\\
 &\ \vdots\\
 r_{n-2}&=q_nr_{n-1}+r_n,\\
 r_{n-1}&=q_{n+1}r_n
\end{aligned}
\]
에서 마지막 영이 아닌 나머지 $r_n$을 일계수로 만든 것이 $\gcd(f,g)$이다.
\end{theorem}

\begin{proofidea}
각 나눗셈 $a=qb+r$에서 $a,b$의 공약수와 $b,r$의 공약수는 같다. 차수가 계속 감소하므로 과정은 유한 번 만에 끝난다.
\end{proofidea}

\begin{theorem}[베주 항등식]
$f,g\in K[X]$에 대해 어떤 $a,b\in K[X]$가 존재하여
\[
  \gcd(f,g)=af+bg
\]
가 성립한다. 특히
\[
  \gcd(f,g)=1
  \quad\Longleftrightarrow\quad
  af+bg=1\text{인 }a,b\in K[X]\text{가 존재한다.}
\]
\end{theorem}

\begin{proof}
$K[X]$가 PID이므로 $(f,g)=(d)$인 일계수 $d$가 존재한다. $d\in(f,g)$이므로 $d=af+bg$이고, $f,g\in(d)$이므로 $d$는 두 다항식의 공약수이다. 모든 공약수는 $af+bg$도 나누므로 $d$가 최대공약수이다.
\end{proof}

\begin{example}
$\Q[X]$에서 $f=X^3-1$, $g=X^2-1$이라 하자. 나눗셈을 하면
\[
 f=Xg+(X-1),\qquad g=(X+1)(X-1).
\]
따라서
\[
  \gcd(f,g)=X-1.
\]
첫 식을 정리하면 $X-1=f-Xg$이므로 베주 표현도 동시에 얻어진다.
\end{example}

\subsection{몫환에서 역원 구하기}

\begin{proposition}
$f,m\in K[X]$이고 $m\ne0$이라 하자. 잉여류 $\overline f\in K[X]/(m)$가 단위원일 필요충분조건은
\[
  \gcd(f,m)=1
\]
이다.
\end{proposition}

\begin{proof}
$af+bm=1$이면 몫환에서 $\overline a\,\overline f=\overline1$이다. 반대로 $\overline f$가 역원을 가지면 어떤 $a$에 대해 $af-1\in(m)$이므로 $af+bm=1$인 $b$가 존재한다.
\end{proof}

\begin{example}
$A=\F_2[X]/(X^3+X+1)$에서 $\overline{X+1}$의 역원을 구하자. 계산하면
\[
 X^3+X+1=(X^2+X)(X+1)+1.
\]
따라서
\[
  1=(X^3+X+1)+(X^2+X)(X+1)
\]
이고, 표수가 $2$이므로
\[
  \overline{X+1}^{\,-1}=\overline{X^2+X}.
\]
\end{example}

\begin{keyidea}
유클리드 알고리즘은 최대공약수만 계산하는 절차가 아니다. 나눗셈을 거꾸로 대입하면 베주 계수를 얻고, 이는 합동식 풀이와 몫환의 역원 계산으로 이어진다.
\end{keyidea}
